% ==========================================================================
% EMPR · Discussion
% 数字全部走 paper/numbers.tex 的宏（由 results/ 生成），本文件不手抄任何数字。
% ==========================================================================
\section{Discussion}
\label{sec:discussion}

\paragraph{From a binary question to a third regime.}
We opened with ``compute or retrieve?'' because it makes falsifiable predictions.
The data reject both endpoints and support a third regime. A pure retrieval story
predicts insensitivity to numerical values and high sensitivity to surface form; a
pure computation story predicts the opposite. We observe \emph{both} sensitivities
at once: rescaled values move \pctChangedSevenBNumeric\% of answers (so the values
are used), and a semantics-preserving paraphrase moves
\pctChangedSevenBWording\% (so surface form is also used). The additional
separation we introduce --- counting only changes that cross the
correct/incorrect boundary --- shows that the paraphrase effect is
\emph{value-level}: just \nbWordingFlips{} of 240 items change correctness, with
discordant pairs roughly balanced. We therefore name the regime
\emph{surface-conditioned computation}: the model computes, its computation depends
on irrelevant phrasing, and its outputs wander within a band that mostly keeps
correctness fixed. Calling this ``reasoning failure'' or ``retrieval'' would both
misdescribe it.

\paragraph{What the evidence jointly supports.}
Five probes converge on one reading. (i) Accuracy on SciBench physics tops out at
\accCotFourteenB\% for the largest model we ran and gains only
\gapFourteenVsSevenCot{} points from a two-fold parameter increase, while
chain-of-thought multiplies accuracy by \ratioThreeB--\ratioFourteenB$\times$ for
the Qwen family --- prompting outweighs scale in this setting. (ii) The same model
reaches \accMMLU\% on multiple-choice MMLU physics but only \accCotSevenB\% when it
must produce the number itself, which rules out missing knowledge: it is present
and retrievable, while numerical execution is not. (iii) The paraphrase
instability described above is model-generic
(\pctChangedThreeBWording\%, \pctChangedFourteenBWording\%,
\pctChangedLlamaEightBWording\% in the other three models) and is value-level
rather than correctness-level. (iv) Restoring clean residual states recovers
behaviour over an early-middle band (layers 4--16 of \nbLayers) that we can
localise as a range but not as a single peak, and the recovery collapses after
layer 18. (v) Separating output execution from correctness shows parseability is
layer-invariant (\pctParseCorrupt\% even under corruption) while correctness is
depth-dependent, and on problems the model cannot solve restoration never exceeds
\pctAccCleanWrongBest\% at any single layer. The failures are capability limits,
not misplaced representations that a layer patch could repair.

\paragraph{Why the error taxonomy is itself uncertain (and why that is a
contribution).}
Our cross-annotator agreement ($\kappa=\kappaAnnotators$, $\kappa_w=\kappaLinearWeighted$
linearly weighted, \pctAgreeAnnotators\% raw agreement over \nbAnnotated{} samples)
is low, and this is not merely a nuisance for our own analysis --- it reveals a
fundamental limitation of the taxonomy itself. The per-class breakdown is
diagnostic: the judges agree \pctAgreeEtwo\% of the time on E2 but only
\pctAgreeEone\% on E1. That is exactly what one expects if the boundary between
``misapplying a physical law'' and ``algebraic miscalculation'' is intrinsically
ambiguous: when a model applies the wrong law, the resulting arithmetic is wrong
too, so the observable text is consistent with either label. We therefore report
only the coarse, annotator-robust statements --- E2 is the modal category under
both judges, E3 is rare, and 7B/14B do not differ under either labelling
($p=\pChiSameFamily$ same-family, $p=\pChiCrossFamily$ cross-family) --- and
deliberately do not quote precise category shares. This suggests that evaluating
reasoning purely through coarse error categories is insufficient: taxonomy-based
analysis inherits an annotator-dependent ceiling that stays invisible unless
agreement is measured. We release our dual-annotated sample (\nbAnnotated{} wrong
answers from two models, each labelled independently by a same-family and a
cross-family judge) as a resource not only for physics reasoning but for research
on LLM-as-a-judge attribution, where the \emph{disagreement} between judges is the
signal rather than noise.

\paragraph{Implications for mitigation.}
Two findings are directly actionable. First, since the dominant failure mode is
arithmetic rather than physical (E2), the highest-leverage interventions are
execution-side --- tool use, code interpreters, or explicit symbolic bookkeeping
--- rather than more physical pre-training data. Second, since
\pctChangedSevenBWording\% of answers move under a semantics-preserving paraphrase,
single-prompt evaluation understates the variance of a model's competence;
reporting paraphrase-consistency alongside accuracy is a cheap and informative
robustness metric. We also tested whether the instability is recoverable at
inference time by adding structural scaffolding to the prompt: at a matched token
budget the change is \scafDeltaPP{} points (McNemar $p=\scafP$), so prompt-level
scaffolding does not substitute for capability, which further argues that the
mitigation has to be execution-side rather than wording-side.

\paragraph{Threats to validity.}
All numbers come from a single GPU class (A800-80G) and a \texttt{transformers}
backend without continuous batching, so wall-clock times are not comparable to a
vLLM deployment, although greedy decoding makes accuracies backend-independent.
Causal tracing covers one open-weight model and the strict-inclusion subset
(\nbCausalValid{} of \nbCausalAttempts{} attempted), so the layer localisation is a
qualitative early-middle pattern rather than a precise per-layer estimate; the
layer-parse analysis rests on \nbCleanCorrect{} correct-baseline problems, which we
accompany with Wilson intervals. OlympiadBench results refer to the
\pctOlympiadScorable\% of problems with numerically gradable references and are used
only for the direction of the scale trend. Error-type shares are annotator
dependent ($\kappa=\kappaAnnotators$) and are reported coarsely. All models are
evaluated zero-shot with a single fixed prompt per condition; prompt sensitivity is
itself one of the phenomena measured, so no per-model prompt tuning is attempted.
